% texte.tex — f173v, vv. 920–958 (conversion initiale des travaux antérieurs)
\vers[n=920]{ΚΡ.}{Kritylla}{οἴμ᾽ ὡς πανοῦργος καὐτὸς εἶναί μοι δοκεῖς}{Ah! toi aussi, que tu m'as l'air d'un fripon}
\vers[n=921]{}{}{καὶ τοῦδέ τις ξύμβουλος. οὐκ ἐτὸς πάλαι}{et d'être son conseiller. Rien d'étonnant que depuis un temps}
\vers[n=922]{}{}{ᾐγυπτιάζετ᾽. ἀλλ᾽ ὅδε μὲν δώσει δίκην.}{vous jouiez les Égyptiens. Mais celui-là en portera la peine.}
\vers[n=923]{}{}{προσέρχεται γὰρ ὁ πρύτανις χὠ τοξότης.}{Car voici venir le prytane et l'archer.}
\vers[n=924]{ΕΥ.}{Euripide}{τουτὶ πονηρόν· ἀλλ᾽ ὑπαποκινητέον.}{Voilà qui ne vaut rien. Il est temps d'esquiver.}
\vers[n=925]{ΜΝ.}{Mnésiloque}{ἐγὼ δ᾽ ὁ κακοδαίμων τί δρῶ;}{Et moi, malheureux, que vais-je faire ?}
\vers[suite]{ΕΥ.}{Euripide}{μέν᾽ ἥσυχος.}{Sois en paix.}
\vers[n=926]{}{}{οὐ γὰρ προδώσω σ᾽ οὐδέποτ᾽, ἤνπερ ἐμπνέω,}{Jamais je ne te trahirai, tant que j'aurai le souffle,}
\vers[n=927]{}{}{ἢν μὴ προλίπωσ᾽ αἱ μυρίαι με μηχαναί.}{à moins que mes mille stratagèmes ne viennent à me faillir.}
\vers[n=928]{ΜΝ.}{Mnésiloque}{αὕτη μὲν ἡ μήρινθος οὐδὲν ἔσπασεν.}{Voilà une amorce qui n'a pas pris.}
\vers[n=929]{ΠΡ.}{Prytane}{ὅδ᾽ ἔσθ᾽ ὁ πανοῦργος ὃν ἔλεγ᾽ ἡμῖν Κλεισθένης;}{Est-ce là le fripon dont nous parlait Clisthène ?}
\vers[n=930]{}{}{οὗτος τί κύπτεις; δῆσον αὐτὸν εἰσάγων}{Holà, pour quoi courbes-tu le front? Archer, rentre avec lui et attache-le}
\vers[n=931]{}{}{ὦ τοξότ᾽ ἐν τῇ σανίδι, κἄπειτ᾽ ἐνθαδὶ}{à l'ais. Puis expose-le ici,}
\vers[n=932]{}{}{στήσας φύλαττε καὶ προσιέναι μηδένα}{et garde-le bien. Ne laisse approcher personne,}
\vers[n=933]{}{}{ἔα πρὸς αὐτόν, ἀλλὰ τὴν μάστιγ᾽ ἔχων}{mais frappe de ton fouet}
\vers[n=934]{}{}{παῖ᾽ ἢν προσίῃ τις.}{quiconque s'avancera.}
\vers[suite]{ΚΡ.}{Kritylla}{νὴ Δί᾽ ὡς νῦν δή γ᾽ ἀνὴρ}{Bien, par Zeus, car tout à l'heure}
\vers[n=935]{}{}{ὀλίγου μ᾽ ἀφείλετ᾽ αὐτὸν ἱστιορράφος.}{un rapetasseur de voiles a failli me l'enlever.}
\vers[n=936]{ΜΝ.}{Mnésiloque}{ὦ πρύτανι πρὸς τῆς δεξιᾶς, ἥνπερ φιλεῖς}{Ô prytane, au nom de cette droite dont tu aimes}
\vers[n=937]{}{}{κοίλην προτείνειν ἀργύριον ἤν τις διδῷ,}{à tendre le creux (À part :) quand on t'offre de l'argent...}
\vers[n=938]{}{}{χάρισαι βραχύ τί μοι καίπερ ἀποθανουμένῳ.}{Accorde-moi une petite faveur, à moi qui suis sur le point de mourir.}
\vers[n=939]{ΠΡ.}{Prytane}{τί σοι χαρίσωμαι;}{Quelle faveur ?}
\vers[suite]{ΜΝ.}{Mnésiloque}{γυμνὸν ἀποδύσαντά με,}{ordonne à l'archer}
\vers[n=940]{}{}{κέλευε πρὸς τῇ σανίδι δεῖν τὸν τοξότην,}{de m'attacher nu face à l'ais.}
\vers[n=941]{}{}{ἵνα μὴ 'ν κροκωτοῖς καὶ μίτραις γέρων ἀνὴρ}{Je crains d'être la risée des crocotes et des mitres}
\vers[n=942]{}{}{γέλωτα παρέχω τοῖς κόραξιν ἑστιῶν.}{avec les corbeaux que je vais régaler.}
\vers[n=943]{ΠΡ.}{Prytane}{ἔχοντα ταῦτ᾽ ἔδοξε τῇ βουλῇ σε δεῖν,}{La boulé a décidé qu'on t'attacherait dans ce costume,}
\vers[n=944]{}{}{ἵνα τοῖς παροῦσι δῆλος ᾖς πανοῦργος ὤν.}{pour témoigner aux personnes présentes que tu es un fripon.}
\vers[n=945]{ΜΝ.}{Mnésiloque}{ἰαππαπαιάξ· ὦ κροκώθ᾽ οἷ᾽ εἴργασαι·}{Aïe, aïe! aïe, aïe! Ô crocote, que de malheurs tu as causés !}
\vers[n=946]{}{}{κοὐκ ἔστ᾽ ἔτ᾽ ἐλπὶς οὐδεμία σωτηρίας.}{Il n'est plus pour moi le moindre espoir de salut.}
\vers[n=947, metre={an⁴}]{ΚΟ.}{Coryphée}{ἄγε νυν ἡμεῖς παίσωμεν ἅπερ νόμος ἐνθάδε ταῖσι γυναιξίν,}{Allons, livrons-nous à nos danses, comme c'est ici la coutume des femmes,}
\vers[n=948]{}{}{παραχώρει ὅταν ὄργια σεμνὰ θεαῖν ἱεραῖς ὥραις ἀνέχωμεν, ἅπερ καὶ}{quand aux dates saintes nous célébrons les augustes mystères des deux déesses.}
\vers[n=949, indent=3em, findent=1.7999999999999998em, metre={an.}]{}{}{Παύσων σέβεται καὶ νηστεύει,}{Pauson aussi les sanctifie. Il jeûne}
\vers[n=950, indent=3em, findent=1.7999999999999998em]{}{}{πολλάκις αὐταῖν ἐκ τῶν ὡρῶν}{et souvent prie avec nous que d'année}
\vers[n=951, indent=3em, findent=1.7999999999999998em]{}{}{ἐς τὰς ὥρας ξυνεπευχόμενος}{en année il puisse longtemps}
\vers[n=952, indent=3em, findent=1.7999999999999998em]{}{}{τοιαῦτα μέλειν θάμ᾽ ἑαυτῷ.}{encore leur témoigner même dévotion.}
\vers[n=953, indent=2em, findent=1.2em, metre={lyr.}]{ΧΟ.}{Chœur}{ὅρμα χώρει·}{Pars, élance-toi}
\vers[n=954, indent=2em, findent=1.2em]{}{}{κοῦφα ποσὶν ἄγ᾽ ἐς κύκλον,}{d'un pied léger, forme le cercle,}
\vers[n=955, indent=2em, findent=1.2em]{}{}{χειρὶ σύναπτε χεῖρα,}{les mains enlacées aux mains,}
\vers[suite]{ΧΟ.}{Chœur}{†ῥυθμὸν χορείας ὕπαγε πᾶσα†·}{marque le rythme de la danse.}
\vers[n=956, indent=2em, findent=1.2em]{}{}{βαῖνε καρπαλίμοιν ποδοῖν.}{Avance de tes pieds alertes.}
\vers[n=957, indent=2em, findent=1.2em]{}{}{ἐπισκοπεῖν δὲ πανταχῇ}{Mais il faut que le chœur soit aux aguets}
\vers[n=958, indent=2em, findent=1.2em]{}{}{κυκλοῦσαν ὄμμα χρὴ χοροῦ κατάστασιν.}{et tourne la vue de toutes parts.}
